\documentclass[11pt,a4paper]{article}

\usepackage[utf8]{inputenc}
\usepackage[T1]{fontenc}
\usepackage[french]{babel}
\usepackage{geometry}
\usepackage{enumitem}
\usepackage{array}
\usepackage{amsmath}
\usepackage{xcolor}
\usepackage{hyperref}
\usepackage{listings}
\usepackage{tikz}
\usetikzlibrary{arrows.meta, positioning, fit}

\geometry{margin=2.2cm}
\hypersetup{colorlinks=true,linkcolor=blue,urlcolor=blue}
\definecolor{softblue}{RGB}{225,238,249}
\definecolor{softgreen}{RGB}{229,242,231}
\definecolor{softorange}{RGB}{252,237,214}
\definecolor{softpurple}{RGB}{239,231,247}
\definecolor{softgray}{RGB}{242,244,246}
\lstset{
  basicstyle=\ttfamily\small,
  columns=fullflexible,
  keepspaces=true,
  frame=single,
  breaklines=true
}

\title{\textbf{Rapport d'analyse}\\
       \large Robot de réconfort}
\author{BOUBENIA Walid \and BARRE-GUETROT Ilan}
\date{}

\begin{document}

\maketitle

\section{Architecture}

L'architecture du projet est organisée autour d'un contrôleur principal et de
composants spécialisés :

\begin{center}
\begin{tikzpicture}[
  node distance=8mm and 12mm,
  box/.style={
    draw=black!65,
    rounded corners=2pt,
    align=center,
    minimum height=10mm,
    text width=32mm,
    font=\small
  },
  input/.style={box, fill=softgray},
  control/.style={box, fill=softblue, text width=48mm, minimum height=12mm},
  service/.style={box, fill=softorange},
  agent/.style={box, fill=softgreen, text width=42mm},
  state/.style={box, fill=softpurple, text width=38mm},
  output/.style={box, fill=softgray},
  arrow/.style={-{Latex[length=2.5mm]}, thick, black!65}
]
  \node[input] (json) {Entrées JSON\\carte, scénario, données};
  \node[control, below=of json] (main)
    {\textbf{agent.py}\\chargement et orchestration};

  \node[service, below left=13mm and 3mm of main] (validation)
    {\textbf{validation.py}\\contrôle de cohérence};
  \node[service, below right=13mm and 3mm of main] (io)
    {\textbf{reconfort\_io.py}\\lecture et simulateur local};

  \node[agent, below=25mm of main] (robot)
    {\textbf{Agent}\\perception, décision et actions};
  \node[state, below=of robot] (memory)
    {\textbf{Mémoire}\\CarteMentale + MemoireRobot};
  \node[state, below=of memory] (planning)
    {\textbf{Planification}\\BFS et sélection du casier};
  \node[output, below=of planning] (trace)
    {\textbf{Trace JSON}\\pas, livraisons, résumé};

  \draw[arrow] (json) -- (main);
  \draw[arrow] (main) -- (validation);
  \draw[arrow] (main) -- (io);
  \draw[arrow] (main) -- (robot);
  \draw[arrow] (io) -- (robot)
    node[midway, right, font=\scriptsize, align=left] {perception\\et casiers};
  \draw[arrow] (validation) |- (robot);
  \draw[arrow] (robot) -- (memory);
  \draw[arrow] (memory) -- (planning);
  \draw[arrow] (planning) -- (trace);
\end{tikzpicture}
\end{center}
\vspace{2mm}
\begin{center}
\small
\textit{Les flèches représentent le flux de contrôle et d'informations.}
\end{center}

Les frontières entre les composants sont volontaires :

\begin{itemize}
  \item \textbf{Perception :} \texttt{PerceptionLocale}, dans
  \href{run:reconfort_io.py}{\texttt{reconfort\_io.py}}, est la seule
  interface qui consulte la grille réelle. Elle expose uniquement les quatre
  cases voisines et le contenu du casier lorsque le robot est devant
  l'armoire. Cela respecte la contrainte de perception locale.

  \item \textbf{Mémoire :} \texttt{CarteMentale} et
  \texttt{MemoireRobot}, dans \texttt{agent.py}, stockent ce que le robot a
  découvert et son état courant. Cette séparation empêche la planification de
  lire directement la carte réelle et permet de tester la mémoire
  indépendamment du simulateur.

  \item \textbf{Décision :} \texttt{Agent.choisir\_action} utilise la
  perception et la mémoire pour choisir \texttt{AVANCER},
  \texttt{CONSULTER}, \texttt{CHERCHER}, \texttt{PRENDRE},
  \texttt{DONNER} ou \texttt{ATTENDRE}. La décision ne modifie pas
  directement la grille simulée.

  \item \textbf{Action :} \texttt{executer\_action} applique l'action à
  l'état interne, appelle l'interface de prise si nécessaire et journalise
  l'état avant l'action. \texttt{Trace} est séparée afin que la journalisation
  n'influence pas la décision.

  \item \textbf{Validation :} \texttt{validation.py} est exécuté avant la
  simulation. Il vérifie la cohérence des cartes, scénarios, dictionnaires et
  casiers, au lieu de disperser ces contrôles dans l'agent.
\end{itemize}

Le flux principal est donc :
\[
\begin{aligned}
\texttt{charger\_entrees}
&\longrightarrow \texttt{valider\_entrees}
\longrightarrow \text{création du simulateur et de l'agent}\\
&\longrightarrow \texttt{percevoir}
\longrightarrow \text{mémoriser}
\longrightarrow \texttt{choisir\_action}\\
&\longrightarrow \texttt{executer\_action}
\longrightarrow \texttt{Trace}.
\end{aligned}
\]

\section{Stratégie de replanification}

Le robot \textbf{replanifie entièrement} son chemin, plutôt que de réparer une
partie du plan. À chaque itération de \texttt{traiter\_demande}, il observe les
quatre voisins, met à jour la carte mentale, puis appelle indirectement
\texttt{plus\_court\_chemin}, qui effectue un BFS depuis la position courante
vers une case adjacente à la cible.

Lorsqu'un mur est découvert, \texttt{memoriser\_perception} inscrit la nature
\texttt{mur} dans \texttt{CarteMentale}. Le BFS suivant exclut alors cette
case et reconstruit un nouveau chemin avec les informations disponibles. Les
cases encore inconnues sont provisoirement considérées comme traversables ;
elles seront corrigées par les perceptions suivantes.

Le coût d'une replanification est
\(\mathcal{O}(H\times L)\) dans le pire cas pour une carte de \(H\times L\)
cases, avec \(\mathcal{O}(H\times L)\) de mémoire pour les parents et les
cases visitées. Cette solution est plus simple et plus sûre qu'une réparation
locale : elle garantit que le chemin utilisé tient compte de tous les murs
connus. Son inconvénient est de recalculer le chemin après chaque action, même
si la modification de la carte est faible. Pour les petites grilles du
projet, ce coût reste acceptable.

\section{Critique de la règle de repli}

La règle imposée en section 6.2 est déterministe, mais elle est discutable :
une émotion demandée est remplacée par une émotion voisine sur la roue, sans
tenir compte de la signification de l'objet. Une personne triste peut donc
recevoir un objet associé au dégoût. La proximité sur la roue de Plutchik ne
garantit pas une proximité de besoin ou de confort. La règle est défendable
comme solution de secours lorsque le casier exact est vide, mais elle ne
devrait pas être interprétée comme une équivalence psychologique.

Une meilleure règle serait :

\begin{enumerate}
  \item chercher d'abord un casier de l'émotion exacte et de l'intensité
  exacte ;
  \item si aucun objet n'est disponible, chercher l'émotion exacte avec une
  intensité voisine ;
  \item ensuite seulement, utiliser une émotion voisine, en privilégiant une
  table de compatibilité sémantique définie pour le projet ;
  \item si aucun objet compatible n'existe, signaler explicitement l'échec
  plutôt que choisir automatiquement une émotion potentiellement inadaptée.
\end{enumerate}

Dans l'implémentation actuelle, la distance émotionnelle est limitée à
\textbf{2 positions} sur la roue. Ainsi, une distance de \(0\) correspond à
l'émotion exacte, une distance de \(1\) à une émotion voisine et une distance
de \(2\) à une émotion située deux positions plus loin. Les émotions à
distance \(3\) ou \(4\) sont ignorées. La distance est circulaire : la
première et la dernière émotion de la liste sont donc voisines. Cette limite
est appliquée dans \texttt{Agent.\_casiers\_candidats} par la condition
\texttt{if distance > 2}.

Il n'existe pas de limite stricte équivalente pour l'intensité : les trois
niveaux restent candidats, mais ils sont classés selon leur écart avec
l'intensité demandée. Le tri privilégie donc d'abord l'émotion exacte, puis
l'intensité la plus proche.

Pour mettre cette règle en œuvre, il faudrait modifier
\texttt{Agent.\_casiers\_candidats} dans \texttt{agent.py} afin de produire des
priorités par niveaux de compatibilité, au lieu d'utiliser uniquement la
distance circulaire. 



\section{Trois situations où le robot se comporte mal}

\subsection{Échec explicite : une émotion inconnue}

Si une demande ne contient aucun mot du dictionnaire,
\texttt{\_lire\_emotion\_message} renvoie \texttt{None} et
\texttt{emotion\_indeterminee} est placé à \texttt{True}. La demande est alors
classée en \texttt{ATTENDRE} et enregistrée comme un échec avec le motif
\texttt{emotion indeterminee}. Le robot ne peut donc pas choisir d'objet ni
effectuer la livraison. Ce comportement est volontaire : sans émotion
identifiée, une réponse automatique serait arbitraire.

Un ancien problème d'état mémoire a été corrigé : le drapeau
\texttt{emotion\_indeterminee} est maintenant remis à \texttt{False} au début de
\texttt{traiter\_demande}. Une demande inconnue ne bloque donc plus les
demandes valides suivantes.

\subsection{Échec explicite : une armoire vide}

Après identification de l'émotion, le robot explore les casiers candidats.
Lorsque tous les casiers compatibles ont été observés comme vides,
\texttt{choisir\_action} renvoie \texttt{ATTENDRE}. La livraison est alors
enregistrée avec le motif \texttt{armoire vide}.

Ce cas est correctement signalé, mais aucune solution de remplacement n'est
proposée lorsque tous les objets pertinents sont indisponibles. Une évolution
possible serait de proposer un objet d'une autre catégorie selon une règle de
compatibilité explicite, ou de transmettre la demande à un autre robot.

\subsection{Échec explicite : un résident inaccessible}

Si aucun chemin vers une case adjacente au résident n'est trouvé dans la carte
mentale, \texttt{\_direction\_vers} renvoie \texttt{None}. Le robot s'arrête
avec l'action \texttt{ATTENDRE} et la trace indique
\texttt{resident inaccessible}. Il ne tente pas de déplacer un mur et ne
produit pas de fausse livraison.

La limite est que \texttt{ATTENDRE} est utilisé comme état terminal : le robot
ne dispose pas d'une stratégie de recherche plus avancée ou de coopération
pour rendre la cible accessible. Pour améliorer ce comportement, il faudrait
ajouter une phase d'exploration contrôlée, ou un motif d'échec distinct
indiquant que la zone a été entièrement explorée.

\paragraph{Autre limite étudiée.}
La règle de repli par distance circulaire peut aussi fournir un objet d'une
émotion voisine peu adaptée, par exemple un objet associé au dégoût pour une
personne triste. Ce point n'est pas un bug par rapport à l'énoncé : le code
applique bien la règle imposée. Il est analysé dans la section précédente
comme une limite de conception.

\end{document}
